% ===================================================================
%  TABELLA nr. 10 — Il mar. Il port.
%  Traduction 2026 d'apres texto/io, controlee sur texto/fr, texto/ca
%  et texto/oc. Consigne : voir texto/rm/10-tabella-01.tex.
%
%  LE TABLEAU 9 FERMAIT UNE MOITIE, CELUI-CI FERME L'AUTRE, ET
%  ENSEMBLE ILS REPONDENT A LA QUESTION POSEE AU TABLEAU 2.
%
%  Le romanche n'a pas de mer. Aucun locuteur n'a jamais vu de sa
%  vallee ce que cette page decrit, et tout le vocabulaire de la page
%  est donc emprunte. LA QUESTION N'EST PAS S'IL EST EMPRUNTE, ELLE
%  EST A QUI.
%
%  Ce n'est pas a l'allemand, qui est pourtant le voisin de tous les
%  jours, celui a qui le romanche a pris « schrinari » l'atelier au
%  tableau 5, « chegls » et « stelzas » les jeux au tableau 2. C'est
%  a l'ITALIEN, et plus exactement au fonds mediterraneen que
%  l'italien a diffuse : « prua », « puppa », « vela », « ancora »,
%  « timun », « port », « mol », « far », « barca », « marinar »,
%  « boia », « flotta ». Ce sont, mot pour mot, les memes que le
%  francais, l'anglais, l'allemand, le russe et le turc ont pris a la
%  meme source.
%
%  ET C'EST LA REPONSE AU « VELO » DU TABLEAU 2. On y avait vu qu'un
%  mot francais pouvait arriver par une route allemande — la langue
%  d'un emprunt ne dit pas le chemin qu'il a pris. Le tableau 10 dit
%  l'inverse et complete : LE CHEMIN NE DIT PAS DAVANTAGE LA LANGUE.
%  Ce qui decide, dans les deux cas, c'est L'ORIGINE DE LA CHOSE. Le
%  velo est arrive par le nord parce qu'il roulait au nord ; le
%  gouvernail est arrive du sud parce qu'il naviguait au sud. La
%  frontiere qu'on partage n'y est pour rien.
%
%  TROIS TABLEAUX, UNE SEULE REGLE :
%    5  — le contact offre, c'est le manque qui prend ;
%    9  — celui qui a l'objet donne le mot ;
%    10 — quand personne sur place n'a l'objet, le mot vient de la ou
%         l'objet a ete fait, et non de la ou passe la frontiere.
%
%  Une seule chose de cette page est romanche sans emprunt, et c'est
%  la seule que la montagne connaisse aussi : « unda », la vague. On
%  la dit d'un torrent comme d'une mer.
% ===================================================================

%%K t10-apar-1 apar
\VUsaut{4.0mm}
\VUtitre{15.0pt}{60}{TABELLA nr. 10}
\VUsaut{3.0mm}
\VUpk{11.4pt}{40}{Il mar. --- Il port.}
\VUsaut{3.0mm}

%%K t10-01-1 p
1. --- La stad, entant che tscherts tschertgan quietezza e frestgezza sin la
muntogna, preferan auters d'ir a la costa. Uschia avain nus fatg l'onn passà,
perquai che nus avain fitg gugent l'\VUgras{Ocean} \textsuperscript{(1)}.

%%K t10-02-1 p
2. --- Quant a mai sent jau in grond plaschair contemplond quella immensa vastezza
d'aua che sa stenda fin a l'\VUgras{orizont} \textsuperscript{(2)}, e vesend partir per ina
lunga \VUgras{traversada} las enormas \VUgras{navs a vapur} \textsuperscript{(3)}, sin las
qualas s'embartgan ils viagiaders che van en America.

%%K t10-03-1 p
3. --- Perquai tscherna mes bab, ch'enconuscha mia inclinaziun, adina, per passentar
las vacanzas, ina plaja fitg quieta, situada sper in da noss gronds \VUgras{ports
militars}.

%%K t10-03-2 p
Durant noss ultim suggiurn è la \VUgras{squadra} vegnida ad ancorar davos l'enorm
\VUgras{dic} \textsuperscript{(4)} che serra e protegia il \VUgras{port militar}
\textsuperscript{(5)}, e jau hai pudì admirar tenor mes giavisch las differentas
\VUgras{navs da guerra}, il pitschen \VUgras{submarin} \textsuperscript{(10)} che sa tuffa e
sparescha sut l'aua ed era l'enorma \VUgras{nav curassada} \textsuperscript{(9)}, che, fin
ussa, aveva quasi mo tema dals attatgas da la \VUgras{torpediera}
\textsuperscript{(8)}.

%%K t10-tit-2 sub
\VUsaut{3.0mm}
\VUpk{10.2pt}{40}{Las pitschnas navs.}
\VUsaut{3.0mm}

%%K t10-04-1 p
4. --- Quella saveva gia \VUgras{abilmain} chattar il punct flaivel da la grossa
curassa da metal per fitgar là il privlus \VUgras{torpedo}, dal qual l'explosiun fa
perir la nav; ma tuttina era ella main da temair ch'il submarin, perquai ch'las
contra-torpedieras al persequiteschan facilmain.

%%K t10-05-1 p
5. --- Jau hai era pudì vesair ils svelts \VUgras{crusaders} \textsuperscript{(6)}
curassads, quasi uschè nunblessabels sco la vaira nav curassada, pliras \VUgras{navs da
transport} \textsuperscript{(12)}, trais u quatter \VUgras{barcas da chanun}
\textsuperscript{(7)} ed a la fin ina \VUgras{fregatta} \textsuperscript{(11)}.
\VUgras{Il spectacul da quella flotta, cur ch'ella manovrescha per levar l'ancora, è il
pli interessant.}

%%K t10-06-1 p
6. --- Tge differenza da construcziun tranter quels gronds mantuns d'atschal e las
elegantas \VUgras{trais-masts} u las graziusas \VUgras{navs a vela} \textsuperscript{(13)},
anc oz duvradas en la flotta da commerzi, ch'jau hai vis entrar en il port, cun tut
lur velas, cur ch'jau spassegiava sin il \VUgras{mol} \textsuperscript{(14)}. Propi han
quellas pers bler da lur avantatgs \VUgras{cur} ch'il \VUgras{vent} era cuntrari. Ellas
han stuì trair giu tut las velas per turnar en il bassin, ed in \VUgras{remorchader}
\textsuperscript{(16)} è stà necessari per las ir a prender e per las manar, sco simplas
\VUgras{barcas da charg} \textsuperscript{(17)}, al cai.

%%K t10-tit-3 sub
\VUsaut{3.0mm}
\VUpk{10.2pt}{40}{Il port da commerzi.}
\VUsaut{3.0mm}

%%K t10-07-1 p
7. --- Quai ch'è interessant en il port dal qual jau discur, è ch'el cuntegna betg mo
in port militar, arrangià artifizialmain cun lavurs gigantescas, ma era in
maravigliusa \VUgras{port da commerzi} situà en la bucca d'in grond flum.

%%K t10-08-1 p
8. --- La lingua da terra che separa ils dus bassins è fortifitgada e defendida
d'ina seria da \VUgras{batterias} e \VUgras{bastiuns} che van enavant sin il mar. Ins ha
da basegn d'ina permissiun speziala per spassegiar sin ils \VUgras{mirs da defensiun}
\textsuperscript{(15)}. Jau l'aveva survegnida facilmain, tras mes aug, ch'è inschigner
dals \VUgras{chantiers navals} \textsuperscript{(18)}, e jau sun ì a visitar las navs ch'ins
constrouescha sut vasts hangars.

%%K t10-09-1 p
9. --- Jau hai schizunt assistì a la lantschada d'ina nav curassada, e jau
m'algord dal mument commovent ch'tut la fulla ha sentì cur ch'la massa enorma,
bandunada a sasezza, ha cumenzà a sbrischar sin ses rom sco ina cuna ed è vegnida a
flottar sin l'aua dal flum.

%%K t10-10-1 p
10. --- Per ir als chantiers navals traversa ins la \VUgras{plazza da manovras}
\textsuperscript{(19)}, nua ch'ils schuldads da la guarnischun vegnan mintga di a sa
exercitar, ins passa sin in \VUgras{puntin} \textsuperscript{(20)} da fier, che collia
l'esplanada da parada cun l'\VUgras{arsenal} \textsuperscript{(21)}.

%%K t10-tit-4 sub
\VUsaut{3.0mm}
\VUpk{10.2pt}{40}{L'arsenal ed ils docs.}
\VUsaut{3.0mm}

%%K t10-11-1 p
11. --- Cun mes aug hai jau visità quella gronda construcziun plaina da
\VUgras{chanuns} \textsuperscript{(22)}, \VUgras{ballas} \textsuperscript{(23)} e \VUgras{granatas}.
Jau hai era vis la \VUgras{fusina da metal} \textsuperscript{(24)}, en la quala vegnan
fatgas las \VUgras{calderas} \textsuperscript{(25)} da las navs a vapur.

%%K t10-12-1 p
12. --- Fitg dasperas èn ils \VUgras{docs} \textsuperscript{(26)} --- docs da reparaziun ---
e las fitg remarcablas \VUgras{docs flottants} \textsuperscript{(27)}, en las qualas jau hai
pudì vesair fitg dasperas, sin ina nav da reparar, l'\VUgras{elica}
\textsuperscript{(28)}, la \VUgras{quiglia} \textsuperscript{(29)} ed il \VUgras{timun}
\textsuperscript{(30)} d'ina barca.

%%K t10-13-1 p
13. --- Il port da commerzi è precis uschè degn da studegiar sco il port militar, e
jau hai passentà bleras uras gapond sin ils cais, nua ch'èn amarradas las \VUgras{navs
a rodas} \textsuperscript{(31)} che van si il flum, e guardond funcziunar las \VUgras{grias
a mast} \textsuperscript{(32)}, che levan sco plimas chargias enormas.

%%K t10-tit-5 sub
\VUsaut{4.0mm}
\VUpk{10.2pt}{40}{Il pilot.}
\VUsaut{3.0mm}

%%K t10-14-1 p
14. --- Jau hai admirà ils \VUgras{transatlantics} \textsuperscript{(34)} che sortivan da la
rada, cun lur \VUgras{bandieras} \textsuperscript{(35)} en l'aria libra, manads d'in abil
\VUgras{pilot} \textsuperscript{(36)} che sa maravigliusamain suandar il chanel marcà cun
\VUgras{boias} \textsuperscript{(40)} e balisas, evitond las \VUgras{gabarras}
\textsuperscript{(41)}, che sa lascha dirigir mo cun fadia cun lur unica vela quadrata.
Jau hai distinguì sin la \VUgras{prua} \textsuperscript{(37)} --- il davant --- las
\VUgras{chabinas} \textsuperscript{(39)} da la segunda classa, e sin la \VUgras{puppa}
\textsuperscript{(38)} --- il davos --- ils saluns per ils passagiers da l'emprima
classa.

%%K t10-15-1 p
15. --- Mintgatant hai jau era assistì al chargiar d'ina \VUgras{barca da charg}
\textsuperscript{(42)}, en la quala vegniva gist amassada la rauba dal \VUgras{magasin}
\textsuperscript{(43)} e da la \VUgras{staziun maritima} \textsuperscript{(44)}, manada là dals
blers trens da la \VUgras{viafierra dal cai} \textsuperscript{(45)}.

%%K t10-tit-6 sub
\VUsaut{3.0mm}
\VUpk{10.2pt}{40}{Il plaschair da vogar.}
\VUsaut{3.0mm}

%%K t10-16-1 p
16. --- Savens hai jau vogà en il \VUgras{doc a flot} \textsuperscript{(53)}, en il qual
vegnan a sa refugiar las \VUgras{barcas} \textsuperscript{(46)}, ed i era ina legria per mai
da manovrar la \VUgras{voga} \textsuperscript{(47)}. Jau sun muntà sin la \VUgras{punt}
\textsuperscript{(48)} d'ina \VUgras{barca a vela} \textsuperscript{(49)}, jau sun rampignà,
cun \VUgras{cordas} \textsuperscript{(52)}, sin il \VUgras{mast} \textsuperscript{(51)} fin en
las \VUgras{antennas} \textsuperscript{(50)}, lura hai jau cuntinuà mia spassegiada
\VUgras{cun
ina iole} \textsuperscript{(54)} tranter pesantas \VUgras{barcas da pestga}
\textsuperscript{(55)}, nua ch'las \VUgras{raits} \textsuperscript{(56)} sitgentan.

%%K t10-17-1 p
17. --- Sin il \VUgras{cai} \textsuperscript{(58)} è passada avant mai la pli differenta
\VUgras{rauba} \textsuperscript{(59)}, cuntegnida en \VUgras{chaschas} \textsuperscript{(60)} u
en \VUgras{ballas} \textsuperscript{(61)}. Ella furmava la \VUgras{chargia}
\textsuperscript{(63)} da las barcunas, e pussantas \VUgras{grias} \textsuperscript{(62)} la
prendevan e la mettevan sin \VUgras{camiuns} \textsuperscript{(64)}, entant ch'ils
vituraders la declaravan e pajavan las taxas en la \VUgras{duana} \textsuperscript{(65)}.

%%K t10-18-1 p
18. --- A la fin, cur ch'jau sun arrivà a la \VUgras{punt che vira} \textsuperscript{(66)}
u a la \VUgras{chalusa} \textsuperscript{(69)}, passond avant la \VUgras{draga}
\textsuperscript{(68)} che nettegia il \VUgras{chanal} \textsuperscript{(67)}, sun jau
debartgà sin il mol sper il \VUgras{semafor} \textsuperscript{(70)}.

%%K t10-19-1 p
19. --- Da l'autra vart da quel mol vegnan propi ad abordar las \VUgras{schalupas}
\textsuperscript{(71)} che manan a terra ils uffiziers da la squadra. I è in spectacul
degn d'admiraziun da vesair ils marinars levar en cadenza lur vogas, entant ch'la
barca, purtada da l'\VUgras{unda} \textsuperscript{(72)}, aborda facilmain la stgala dal
debartgadi. Sin quest lieu po ins meglier observar la \VUgras{marea} ed il muvament dal
\VUgras{flus} e dal \VUgras{reflus} sin la \VUgras{plaja} \textsuperscript{(73)}.

%%K t10-tit-7 sub
\VUsaut{3.0mm}
\VUpk{10.2pt}{40}{Il club.}
\VUsaut{3.0mm}

%%K t10-20-1 p
20. --- Per turnar a chasa hai jau duvrà il \VUgras{tram electric} \textsuperscript{(74)},
dal qual la \VUgras{fermada} \textsuperscript{(75)} è sper il \VUgras{pal} mess avant il club
dals uffiziers.

%%K t10-20-2 p
Dal \VUgras{balcun} \textsuperscript{(76)} dal club, ornà cun \VUgras{ancoras}
\textsuperscript{(77)}, chanuns e ballas, hai jau savens vis ina splendida radunanza
d'uffiziers da mar: \VUgras{lieutenants} \textsuperscript{(78)} cun lur spada,
\VUgras{capitanis d'artillaria} \textsuperscript{(79)} e d'\VUgras{infantaria}
\textsuperscript{(80)} cun lur \VUgras{sabla}. Davos els era in \VUgras{marinar}
\textsuperscript{(81)} ch'als purtava in \VUgras{cannocchial}, cur ch'els vulevan
distinguer quai che capitava dalunsch.

%%K t10-21-1 p
21. --- Jau hai era vis giuvens \VUgras{sutlieutenants} \textsuperscript{(82)} che
retschaivevan las instrucziuns da lur \VUgras{commandant} \textsuperscript{(83)} u da
l'\VUgras{admiral} \textsuperscript{(84)}, entant ch'in \VUgras{quartiermaister}
\textsuperscript{(85)} spetgava per las purtar a la nav.

%%K t10-22-1 p
22. --- Da quel balcun è la vista splendida: ins domina tut la plaja cun sias
\VUgras{tendas} \textsuperscript{(86)} e \VUgras{chabinas} \textsuperscript{(87)}, ed ins
vesa, a l'ura dal bogn, ils \VUgras{bogniders} \textsuperscript{(88)} che petulan
allegramain avant il \VUgras{casino} \textsuperscript{(89)}.

%%K t10-23-1 p
23. --- A la fin da la plaja furma la costa ina pitschna \VUgras{penisla}
\textsuperscript{(91)} \VUgras{greppusa}, nua ch'las \VUgras{undas grondas}
\textsuperscript{(92)} vegnan a sa rumper e sin la quala è construì in \VUgras{far}
\textsuperscript{(93)}, che mussa la notg la via als navigants. Quella penisla è colliada
cun la terra mo tras in \VUgras{istmus} fitg stretg \textsuperscript{(90)}.

%%K t10-tit-8 sub
\VUsaut{3.0mm}
\VUpk{10.2pt}{40}{Vista generala da las costas.}
\VUsaut{3.0mm}

%%K t10-24-1 p
24. --- Il \VUgras{grep da riva} \textsuperscript{(94)} sa stenda, sfess dal mar, che
dissegna en el cun capricci ina seria da \VUgras{bajas} \textsuperscript{(95)} e da
\VUgras{promontoris} (chaus), che al fan fitg pittoresc.

%%K t10-24-2 p
Sin il \VUgras{plateau} \textsuperscript{(96)}, che domina la \VUgras{stretga da mar}
\textsuperscript{(98)}, è in \VUgras{fort} \textsuperscript{(97)} fitg impurtant ed intginas
«villas».

%%K t10-25-1 p
25. --- Quella stretga separa dal continent ina \VUgras{insla} \textsuperscript{(99)} fitg
privlusa, circumdada da \VUgras{scogls} \textsuperscript{(100)} e \VUgras{brisants}, nua che
barcas ed era grondas navs sa cluppan mintgatant. Malgrà la prontezza dals socurs e
l'arrivada svelta dals \VUgras{battels da salvament}, èn quels \VUgras{naufragis}
terrorisants e, durant las \VUgras{tempestas} equinozialas, èn pliras navs peridas cun
umens e rauba.

%%K t10-25-2 p
Malgrà quella vischinanza è \VUgras{il port} fitg frequentà dals marinars.

\VUsaut{6.0mm}
\VUfilet{22mm}
